% Part: propositional-logic
% Chapter: syntax-and-semantics
% Section: valuations-sat

\documentclass[../../../include/open-logic-section]{subfiles}

\begin{document}

\olfileid[fr]{pl}{syn}{val}

\olsection{\usetoken{P}{valuation} et satisfaction}

\begin{defn}[!!^{valuation}s]
Soit $\{\True, \False\}$ l'ensemble des deux valeurs de vérité,
« vrai » et « faux ». Une \emph{!!{valuation}} pour $\Lang{L_0}$
est une fonction~$\pAssign{v}$ qui attribue à chaque
!!{propositional variable} du langage soit $\True$, soit $\False$ :
$\pAssign{v} \colon \PVar \to \{\True, \False\}$.
\end{defn}

\begin{defn}
  Étant donnée une !!{valuation}~$\pAssign{v}$, on définit
  inductivement la fonction d'évaluation
  $\pValue{v} \colon \Frm[L_0] \to \{\True, \False\}$ par :
  \begin{align*}
    \iftag{prvFalse}{\pValue{v}(\lfalse) & = \False; \\}{}
    \iftag{prvTrue}{\pValue{v}(\ltrue) & = \True; \\}{}
    \pValue{v}(\Obj p_n) & = \pAssign{v}(\Obj p_n); \\
    \iftag{prvNot}{\pValue{v}(\lnot !A) & = \begin{cases}
        \True & \text{si } \pValue{v}(!A) = \False;\\
        \False & \text{sinon.}
      \end{cases} \\ }{}
    \iftag{prvAnd}{\pValue{v}(!A \land !B) & = \begin{cases}
        \True &
        \text{si $\pValue{v}(!A) = \True$ et $\pValue{v}(!B) = \True$;}\\
        \False &
        \text{si $\pValue{v}(!A) = \False$ ou $\pValue{v}(!B) = \False$}.
      \end{cases}\\}{}
    \iftag{prvOr}{\pValue{v}(!A \lor !B) & = \begin{cases}
        \True &
        \text{si $\pValue{v}(!A) = \True$ ou $\pValue{v}(!B) = \True$;}\\
        \False &
        \text{si $\pValue{v}(!A) = \False$ et $\pValue{v}(!B) = \False$}.
      \end{cases}\\}{}
    \iftag{prvIf}{\pValue{v}(!A \lif !B) & = \begin{cases}
        \True &
        \text{si $\pValue{v}(!A) = \False$ ou $\pValue{v}(!B) = \True$;}\\
        \False &
        \text{si $\pValue{v}(!A) = \True$ et $\pValue{v}(!B) = \False$}.
      \end{cases}\\}{}
    \iftag{prvIff}{\pValue{v}(!A \liff !B) & = \begin{cases}
        \True &
        \text{si $\pValue{v}(!A) = \pValue{v}(!B)$;}\\
        \False &
        \text{si $\pValue{v}(!A) \neq \pValue{v}(!B)$}.
      \end{cases}}{}
\end{align*}
\end{defn}

\begin{explain}
Ces clauses correspondent aux tables de vérité suivantes :
\begin{center}
\iftag{prvNot}{
  \begin{tabular}{|c||c|} \hline
    $!A$ & $ \lnot !A$ \\
    \hline \hline
    $\True$ & $\False$ \\
    $\False$ & $\True$ \\
    \hline
  \end{tabular}
}{}
\iftag{prvAnd}{
  \begin{tabular}{|cc||c|} \hline
    $!A$ & $!B$ & $!A \land !B$ \\
    \hline \hline
    $\True$ & $\True$ & $\True$ \\
    $\True$ & $\False$ & $\False$ \\
    $\False$ & $\True$ & $\False$ \\
    $\False$ & $\False$ & $\False$ \\
    \hline
  \end{tabular}
}{}
\iftag{prvOr}{
  \begin{tabular}{|cc||c|} \hline
    $!A$ & $!B$ & $!A \lor !B$ \\
    \hline \hline
    $\True$ & $\True$ & $\True$ \\
    $\True$ & $\False$ & $\True$ \\
    $\False$ & $\True$ & $\True$ \\
    $\False$ & $\False$ & $\False$ \\
    \hline
  \end{tabular}
}{}%
\iftag{notprvNot,notprvAnd,notprvOr,notprvIf}{}{\\[1em]}
\iftag{prvIf}{
  \begin{tabular}{|cc||c|} \hline
    $!A$ & $!B$ & $!A \lif !B$ \\
    \hline \hline
    $\True$ & $\True$ & $\True$ \\
    $\True$ & $\False$ & $\False$ \\
    $\False$ & $\True$ & $\True$ \\
    $\False$ & $\False$ & $\True$ \\
    \hline
  \end{tabular}
}{}
\iftag{prvIff}{
  \begin{tabular}{|cc||c|} \hline
    $!A$ & $!B$ & $!A \liff !B$ \\
    \hline \hline
    $\True$ & $\True$ & $\True$ \\
    $\True$ & $\False$ & $\False$ \\
    $\False$ & $\True$ & $\False$ \\
    $\False$ & $\False$ & $\True$ \\
    \hline
  \end{tabular}
}{}
\end{center}
\end{explain}

\begin{prob}
On ajoute à $\Lang{L_0}$ un connecteur ternaire $\diamondsuit$ dont
l'évaluation est donnée par
\begin{gather*}
  \pValue{v}(\diamondsuit( !A , !B , !C ) ) = \begin{cases}
    \pValue{v}( !B ) &
    \text{si $\pValue{v}(!A) = \True$;}\\
    \pValue{v}( !C ) &
    \text{si $\pValue{v}(!A) = \False $}.
  \end{cases}
\end{gather*}
Écrire la table de vérité de ce connecteur.
\end{prob}

\begin{thm}[Détermination locale]
\ollabel{thm:LocalDetermination}
Supposons que les !!{valuation}s $\pAssign{v_1}$ et $\pAssign{v_2}$
coïncident sur les !!{propositional variable}s qui figurent dans
une !!{formula}~$!A$ : autrement dit,
$\pAssign{v_1}(\Obj p_n) = \pAssign{v_2}(\Obj p_n)$ dès que
$\Obj p_n$ figure dans~$!A$. Alors $\pValue{v_1}$ et $\pValue{v_2}$
coïncident aussi sur~$!A$ :
$\pValue{v_1}(!A) = \pValue{v_2}(!A)$.
\end{thm}

\begin{proof}
Par induction sur $!A$.
\end{proof}

\begin{defn}[Satisfaction]
\ollabel{defn:satisfaction}
On peut définir inductivement la \emph{satisfaction d'une
!!{formula}~$!A$ par une !!{valuation}~$\pAssign{v}$},
notée $\pSat{v}{!A}$, de la façon suivante.
(On écrit $\pSat/{v}{!A}$ pour « non $\pSat{v}{!A}$ ».)
\begin{enumerate}
\tagitem{prvFalse}{%
  \indcase{!A}{\lfalse}{$\pSat/{v}{\indfrm}$.}}{}

\tagitem{prvTrue}{%
  \indcase{!A}{\ltrue}{$\pSat{v}{\indfrm}$.}}{}

\item \indcase{!A}{\Obj p_i}{$\pSat{v}{\indfrm}$
  si et seulement si $\pAssign{v}(\Obj p_i) = \True$.}

\tagitem{prvNot}{%
  \indcase{!A}{\lnot !B}{$\pSat{v}{\indfrm}$ si et seulement si
    $\pSat/{v}{!B}$.}}{}

\tagitem{prvAnd}{%
  \indcase{!A}{(!B \land !C)}{$\pSat{v}{\indfrm}$ si et seulement si
    $\pSat{v}{!B}$ et $\pSat{v}{!C}$.}}{}

\tagitem{prvOr}{%
  \indcase{!A}{(!B \lor !C)}{$\pSat{v}{\indfrm}$ si et seulement si
    $\pSat{v}{!B}$ ou $\pSat{v}{!C}$ (éventuellement les deux).}}{}

\tagitem{prvIf}{%
  \indcase{!A}{(!B \lif !C)}{$\pSat{v}{\indfrm}$ si et seulement si
    $\pSat/{v}{!B}$ ou $\pSat{v}{!C}$ (éventuellement les deux).}}{}

\tagitem{prvIff}{%
  \indcase{!A}{(!B \liff !C)}{$\pSat{v}{\indfrm}$ si et seulement si
    $\pSat{v}{!B}$ et $\pSat{v}{!C}$ sont tous deux vrais, ou tous
    deux faux.}}{}
\end{enumerate}
Si $\Gamma$ est un ensemble de !!{formula}s, $\pSat{v}{\Gamma}$
si et seulement si $\pSat{v}{!A}$ pour toute~$!A \in \Gamma$.
\end{defn}

\begin{prop}\ollabel{prop:sat-value}
  $\pSat{v}{!A}$ si et seulement si $\pValue{v}(!A) = \True$.
\end{prop}

\begin{proof}
  Par induction sur~$!A$.
\end{proof}

\begin{prob}
  Démontrer la \olref[pl][syn][val]{prop:sat-value}.
\end{prob}

\end{document}
